\section{Objetivo del Informe}

El objetivo central del presente informe es evaluar los resultados obtenidos tras la implementación del sistema Odoo en PERÚ FLORES TRAVEL AGENCY TOURS OPERADOR E.I.R.L., analizando el impacto generado en sus procesos comerciales y operativos. La implementación tuvo como finalidad transformar el modelo operativo de la agencia mediante la adopción de una plataforma integrada que permita centralizar la información de pasajeros y proveedores, automatizar la generación de paquetes turísticos personalizados y asegurar la trazabilidad integral de cada etapa del viaje.

Asimismo, se buscó dotar a la organización de una infraestructura tecnológica escalable que optimice los tiempos de respuesta comercial ---reduciendo el tiempo de cotización de 45-60 minutos a menos de 15 minutos---, elimine los errores operativos manuales y fortalezca la capacidad de toma de decisiones basada en datos reales. En este contexto, el informe presenta una evaluación estructurada de los beneficios alcanzados, el nivel de adopción del sistema por parte del personal y las mejoras sustanciales en la eficiencia operativa, evidenciando el valor generado por la digitalización en el desempeño global de la empresa.
